\documentclass[11pt]{article}
\usepackage[a4paper,margin=2.5cm]{geometry}
\usepackage[T1]{fontenc}
\usepackage{lmodern}
\usepackage{graphicx}
\usepackage{url}
% Optional packages (texlive-latex-extra); minimal fallbacks keep the draft building without them.
\IfFileExists{microtype.sty}{\usepackage{microtype}}{}
\IfFileExists{booktabs.sty}{\usepackage{booktabs}}{%
  \newcommand{\toprule}{\hline}\newcommand{\midrule}{\hline}\newcommand{\bottomrule}{\hline}%
  \def\cmidrule(##1)##2{\cline{##2}}}
\IfFileExists{xcolor.sty}{\usepackage{xcolor}}{\usepackage{color}}
\usepackage[hidelinks]{hyperref}

% Generated by docs/bench/plot.py from the raw results. Do not edit.
\newcommand{\layersHiCpu}{27.85}
\newcommand{\layersHiLv}{3.01}
\newcommand{\layersHiOff}{2.08}
\newcommand{\layersHiOvl}{1.41}
\newcommand{\layersHiPf}{4.50}
\newcommand{\layersHiRem}{9.12}
\newcommand{\layersHiRm}{2.32}
\newcommand{\layersLoCpu}{0.96}
\newcommand{\layersLoLv}{2.49}
\newcommand{\layersLoOff}{1.46}
\newcommand{\layersLoOvl}{1.38}
\newcommand{\layersLoPf}{3.75}
\newcommand{\layersLoRem}{0.51}
\newcommand{\layersLoRm}{1.61}
\newcommand{\llmAccDecode}{1.17}
\newcommand{\llmAccPrefill}{837}
\newcommand{\llmCpuDecode}{3.63}
\newcommand{\llmCpuPrefill}{25}
\newcommand{\llmLvDecode}{0.58}
\newcommand{\llmLvPrefill}{499}
\newcommand{\llmWeightsGiB}{7.5}
\newcommand{\pcieDtoH}{6.3}
\newcommand{\pcieHtoD}{6.1}
\newcommand{\sweepFitGBps}{304}
\newcommand{\sweepOversubGBps}{8.5}
\newcommand{\usableVram}{4.89}
\newcommand{\yueAccDecode}{1.9}
\newcommand{\yueAccNar}{1.32}
\newcommand{\yueAccPrefill}{607}
\newcommand{\yueCpuDecode}{6.8}
\newcommand{\yueCpuNar}{52.2}
\newcommand{\yueLvDecode}{43.1}
\newcommand{\yueLvDecodeRange}{41.0--46.0}
\newcommand{\yueLvNar}{1.90}
\newcommand{\yueLvOfOracle}{95}
\newcommand{\yueLvOverAcc}{23}
\newcommand{\yueLvPrefill}{611}
\newcommand{\yueNarSlowdown}{1.4}
\newcommand{\yueOracleDecode}{45.3}
\newcommand{\yueOraclePrefill}{1694}
\newcommand{\yueRepeats}{3}

% Fallbacks so the draft builds before every benchmark has run.
\providecommand{\na}{\textbf{??}}
\newcommand{\num}[1]{\ifcsname #1\endcsname\csname #1\endcsname\else\na\fi}
\newcommand{\fig}[2]{\IfFileExists{figures/#2.pdf}{\includegraphics[width=#1]{figures/#2.pdf}}{\fbox{\na\ figure #2 pending}}}
\newcommand{\todo}[1]{\textcolor{red}{[TODO: #1]}}


\title{lowvram: Transparent GPU Memory Oversubscription\\for Unmodified CUDA Programs}
\author{Henrique Soares\\\small\texttt{henriquegogo@gmail.com}}
\date{October 2026 --- draft}

\begin{document}
\maketitle

\begin{abstract}
Programs written for a GPU fail with an out-of-memory error as soon as their allocations exceed
device memory, even when the host has plenty of free RAM. The established remedies---offloading
engines, automatic device placement, custom allocators---all require changes to the program or
support in its framework. We present \emph{lowvram}, a 140-line C library that runs unmodified
CUDA programs beyond device memory. It serves every \texttt{cuMemAlloc} with
\texttt{cuMemAllocManaged}, so the NVIDIA driver's Unified Memory pages data between GPU and host on
demand, and it enlarges the reported device memory so that self-sizing programs use the extra
capacity. One ELF file is both the preload library and its launcher (\texttt{lowvram python
script.py}). On a 6\,GiB laptop GPU, lowvram costs nothing measurable when a program fits. On a real
workload---the 7.3\,GB YuE2-3B music model, which does not load on this GPU---it generates
tokens at \num{yueLvDecode} tokens/s, \num{yueLvOfOracle}\% of the speed of a hand-pruned copy of
the model that fits in VRAM and \num{yueLvOverAcc}$\times$ faster than Hugging Face Accelerate's
offloading, because demand paging keeps resident exactly the half of the weights that token
generation uses. When every step touches more data than fits, paging loses: it is
\num{yueNarSlowdown}$\times$ slower than layer-wise offloading for YuE2's flow-matching step, and on
a dense 4B language model it generates \num{llmLvDecode} tokens per second against
\num{llmAccDecode} for Accelerate and \num{llmCpuDecode} on the CPU. A single driver hint that only
the program can give recovers most of that gap. The deciding property is not the size of the
program but whether the data a phase reuses fits in VRAM.
\end{abstract}

% ---------------------------------------------------------------------------
\section{Introduction}

Consumer and laptop GPUs commonly have 4--12\,GiB of memory, while the models people want to
run---language models, music, image and video generators---often need more. The failure is abrupt:
the framework's allocator asks the driver for device memory, the driver refuses, and the program
aborts with \texttt{CUDA out of memory}, although the machine may have several times that amount
of idle system RAM.

The established remedies require intervention in the program. Frameworks offer offloading engines
\cite{zero-offload,zero-infinity,flexgen} and automatic device placement \cite{accelerate} that
splits a model between GPU and CPU; research systems swap tensors under the control of a planner
\cite{vdnn,swapadvisor,capuchin}; inference engines run part of a model on the CPU
\cite{llamacpp,powerinfer,ktransformers}. Each works for the programs written against it, but a user
holding an arbitrary research repository cannot easily apply any of them.

CUDA Unified Memory \cite{cuda-guide} offers a different mechanism. Memory allocated with
\texttt{cudaMallocManaged} is addressable from CPU and GPU, and on Pascal and later GPUs under Linux
it can exceed device memory: the driver migrates pages to the GPU on fault and evicts the least
recently used ones when device memory is full \cite{zheng2016,ganguly2019}. Programs must opt in,
however, and mainstream frameworks do not.

lowvram closes that gap without touching the program. It is a shared library, loaded with
\texttt{LD\_PRELOAD}, that turns ordinary device allocations into managed ones at the CUDA driver
API. Our contributions are:
\begin{itemize}
  \item A minimal interposition design (\S\ref{sec:design}) that reaches the CUDA runtime, which
        resolves driver functions privately through \texttt{dlsym} and \texttt{cuGetProcAddress}
        and so escapes plain symbol interposition.
  \item A single-file packaging technique: one ELF object that is both the preload library and
        the launcher.
  \item A measurement study on a 6\,GiB laptop GPU (\S\ref{sec:method}--\S\ref{sec:results}) that
        compares transparent paging with manual streaming, framework offloading, CPU execution and
        an oracle, on synthetic kernels and on two real models---a 7.3\,GB music model with two
        execution paths and a dense 4B language model---and tests whether driver hints close the
        gap to hand-written offloading.
  \item A characterization of where transparent oversubscription is valuable and where it is
        not (\S\ref{sec:where}).
\end{itemize}

% ---------------------------------------------------------------------------
\section{Background: four ways to run beyond VRAM}
\label{sec:background}

A device allocation (\texttt{cudaMalloc}) is backed by physical GPU memory for its whole
lifetime; the driver never pages it out, so there is no fallback when memory runs out. Data that
does not fit must live in host RAM, and there are four ways to use it (Table~\ref{tab:strategies}).
The bandwidths that bound them on our test machine are 300\,GB/s for VRAM, roughly 6\,GB/s for
the PCIe link (\S\ref{sec:method}) and tens of GB/s for host RAM.

\begin{table}[h]
\centering\small
\begin{tabular}{p{3.3cm}p{5.2cm}p{5.2cm}}
\toprule
Strategy & What happens to data that does not fit & Examples \\
\midrule
CPU execution & stays in RAM; the CPU computes on it & PyTorch on CPU, llama.cpp \texttt{-ngl 0} \\
Hybrid split & stays in RAM; the CPU computes those layers or experts & llama.cpp \texttt{-ngl N}, KTransformers \\
Explicit offload & copied to the GPU before use, by a plan & Accelerate, ZeRO-Offload, FlexGen \\
Demand paging & migrated to the GPU on access, by the driver & Unified Memory; lowvram \\
\bottomrule
\end{tabular}
\caption{Strategies for working sets larger than device memory. Only demand paging needs no
knowledge of the program.}
\label{tab:strategies}
\end{table}

The first three need knowledge of the computation: which tensors exist, when they are used,
where to compute them. Demand paging needs none; its cost is that transfers start only after a
fault, in small granules, and that eviction follows recency rather than future use.

% ---------------------------------------------------------------------------
\section{Design}
\label{sec:design}

\subsection{Allocation}
PyTorch's caching allocator \cite{pytorch}, like most CUDA programs, obtains device memory through
\texttt{cudaMalloc}, which calls the driver's \texttt{cuMemAlloc\_v2}. lowvram replaces that
function:
{\small
\begin{verbatim}
CUresult cuMemAlloc_v2(CUdeviceptr *dptr, size_t size) {
    init();
    return real_alloc_managed(dptr, size, CU_MEM_ATTACH_GLOBAL);
}
\end{verbatim}
}
The returned pointer is a valid device pointer, so kernels and libraries such as cuBLAS and cuDNN
use it unchanged, and \texttt{cuMemFree} releases it without a hook. While the working set fits,
pages stay resident after their first touch and cost the same as device memory. Beyond that, the
driver evicts pages to host memory and faults them back in on access: an out-of-memory failure
becomes a slowdown.

\subsection{Reported capacity}
Some programs size themselves from \texttt{cudaMemGetInfo}: automatic device placement, for
instance, puts on the CPU whatever does not fit in the reported free memory. lowvram therefore adds
half of physical RAM to the free and total values returned by \texttt{cuMemGetInfo\_v2} and
\texttt{cuDeviceTotalMem\_v2}. The addition is constant: programs that read free memory once
behave correctly, but one that allocates until reported free memory is exhausted can exceed
physical RAM.

\subsection{Reaching the CUDA runtime}
Symbol interposition alone does not reach most CUDA programs. The CUDA runtime does not link
against \texttt{libcuda}: it opens the driver with \texttt{dlopen}, obtains
\texttt{cuGetProcAddress} through \texttt{dlsym}, and resolves every other driver function through
it. lowvram therefore also defines \texttt{dlsym}, \texttt{cuGetProcAddress} and
\texttt{cuGetProcAddress\_v2}. Each calls the real implementation and substitutes a hook when the
requested name is one lowvram replaces. The two lookup paths use different names---the exported
symbol (\texttt{cuMemAlloc\_v2}) for \texttt{dlsym}, the base name (\texttt{cuMemAlloc}) for
\texttt{cuGetProcAddress}---so the substitution table accepts both. The real \texttt{dlsym} is
obtained with \texttt{dlvsym(RTLD\_NEXT, "dlsym", ...)} to avoid recursion. Initialization is lazy
and happens once, on the first hooked call, because other libraries may call \texttt{dlsym} before
lowvram's constructors run; without an NVIDIA driver the library is inert, which matters because
the preload also reaches launchers and shells.

\subsection{Single-file launcher}
lowvram is a shared object that is also executable, as glibc's own \texttt{libc.so.6} is: a
\texttt{.interp} section names the dynamic loader, and a linker-specified entry point
(\texttt{-Wl,-e,lowvram\_main}) runs when the file is executed. The entry point reads its arguments
from \texttt{/proc/self/cmdline} and its own path from \texttt{/proc/self/exe}, prepends that path
to \texttt{LD\_PRELOAD} and \texttt{exec}s the command. Two pitfalls: without C start-up code the
entry point receives no \texttt{argc}/\texttt{argv} and the object's constructors do not run, so
the arguments must come from \texttt{/proc}; and the entry point is reached without a return
address, so it must realign the stack (\texttt{force\_align\_arg\_pointer}) and never return.
Building a position-independent executable and preloading it instead is not possible: glibc 2.30
and later refuse to load PIE executables as libraries. Because \texttt{LD\_PRELOAD} is inherited,
wrappers and children are covered: \texttt{lowvram uv run script.py} reaches the Python interpreter
that \texttt{uv} starts.

\subsection{Scope}
Allocations that bypass \texttt{cuMemAlloc} are not converted: stream-ordered allocation
(\texttt{cudaMallocAsync}) and the virtual memory management API (\texttt{cuMemCreate}/
\texttt{cuMemMap}), which PyTorch uses when \texttt{expandable\_segments} is enabled. These remain
limited to device memory. The implementation targets Linux with glibc on x86-64.
\todo{AArch64: fix the dlvsym fallback (GLIBC\_2.17) and test before claiming it.}

% ---------------------------------------------------------------------------
\section{Methodology}
\label{sec:method}

\paragraph{Machine.} NVIDIA GeForce RTX 3060 Laptop GPU (GA106, 6\,GiB, compute capability 8.6),
Intel Core i7-11800H (8 cores, 16 threads), 31\,GiB RAM, Linux 7.0, NVIDIA driver 595.91.07,
glibc 2.43, PyTorch 2.14.1 with CUDA 13.0, Transformers 4.57.6. The GPU also drives the
display, which holds about 0.7\,GiB, so \num{usableVram}\,GiB were free to the benchmarks. Its link
is PCIe 4.0 x16-capable but trained at x8 by the laptop; pinned host-to-device copies reach
\num{pcieHtoD}\,GB/s and device-to-host \num{pcieDtoH}\,GB/s. This bandwidth, not the GPU,
bounds every strategy that moves data.

\paragraph{Procedure.} Every configuration runs in a fresh process, with and without lowvram
preloaded. Synthetic benchmarks report the median of five timed repetitions after a warm-up
(twenty for GEMM, three or two for slow CPU runs); the YuE2 benchmarks repeat each configuration
in \num{yueRepeats} separate processes and report the median. An out-of-memory error is recorded as
a result. The benchmark scripts, the raw JSON results and the script that turns them into every
figure, table and number in this paper are in \texttt{docs/bench} of the source repository.

\paragraph{Synthetic workloads.}
(i) \emph{GEMM}: an $8192\times8192$ matrix product in FP16 and FP32.
(ii) \emph{Sweep}: a working set of 256\,MiB tensors, each incremented in place once per pass in a
fixed order; we report effective bandwidth (bytes read plus written per second).
(iii) \emph{Layer stack}: a chain of $8192\times8192$ FP16 matrix products standing in for
transformer inference, with 2--12\,GiB of weights and batch 16 (memory-bound) or 2048
(compute-bound). It runs in seven modes: \emph{gpu} (weights allocated on the GPU; unmodified
PyTorch or lowvram), \emph{manual streaming} (weights in pinned RAM, each layer copied to the GPU
just before use), \emph{overlapped streaming} (the copy of layer $i{+}1$ overlaps the product of
layer $i$ on a second stream, the standard double-buffering scheme), \emph{CPU} (FP32, because
the CPU lacks fast FP16 arithmetic), and three \emph{hint} modes that run under lowvram and add one
driver hint from the program: \texttt{READ\_MOSTLY} on all weights (evicting a read-duplicated
page drops it instead of copying it back), \emph{prefetch} (\texttt{cuMemPrefetchAsync} of the
next two layers on a side stream), and \emph{remote map} (the layers that do not fit a 3.5\,GiB
budget are pinned to host RAM and mapped for direct access, so they never migrate). The hint modes
measure how much an informed program could recover; lowvram itself issues no hints.

\paragraph{Real workload: YuE2-3B.} YuE2 \cite{yue2} is a music generation model whose
3.6 billion BF16 parameters occupy 7.26\,GB (6.76\,GiB)---more than this GPU has. It is an AR--NAR
``mixture of transformers'': each of its 28 layers holds two attention blocks and two MLPs, one
pair for autoregressive token generation (AR) and one for a flow-matching decoder (NAR). Token
generation uses only the AR pair, the embeddings and the output head, about 3.4\,GiB; a
flow-matching step runs both pairs over the whole sequence. We measure both phases with the model's
own code through Transformers: AR prefill of 512 tokens and greedy decoding of 96 tokens with a KV
cache (excluding the first 16 from the median), and NAR velocity evaluations over a 1,538-token
sequence (512 context tokens, 1,024 latent frames). Inputs are synthetic token ids and latents of
the real shapes, since timing depends on shapes, not values. Modes: unmodified PyTorch;
lowvram; Accelerate \cite{accelerate} with \texttt{device\_map="auto"} and a 4\,GiB weight budget,
the standard way to run an oversized Hugging Face model; CPU in BF16; and, for token generation, an
\emph{oracle} that deletes the NAR half before moving the model to the GPU, so that the AR half
runs from ordinary device memory. The oracle requires knowing the model's structure; it is the
speed lowvram would reach with perfect placement.

\paragraph{Real workload: a dense language model.} YuE2's two paths favor demand paging, so we add
the general case: Qwen3-4B, a dense decoder whose \num{llmWeightsGiB}\,GiB of BF16 weights are all
read for every token. We measure prompt processing of 1,024 tokens (compute-bound) and greedy
generation of 48 tokens (memory-bound) through Transformers, in the same modes as YuE2 except the
oracle, which does not exist for a model with no idle half.

% ---------------------------------------------------------------------------
\section{Results}
\label{sec:results}

\subsection{No cost when the working set fits}

\begin{table}[h]
\centering\small
\begin{tabular}{lrrr}
\toprule
Workload & Unmodified & lowvram & Change \\
\midrule
GEMM FP16 (TFLOP/s) & 26.29 & 26.29 & -0.0\% \\
GEMM FP32 (TFLOP/s) & 8.82 & 8.79 & -0.3\% \\
Sweep 1 GiB (GB/s) & 309 & 308 & -0.2\% \\
Sweep 4 GiB (GB/s) & 301 & 303 & +0.7\% \\
Layers 2 GiB, batch 16 (ms) & 7.6 & 7.4 & -2.3\% \\
Layers 4 GiB, batch 16 (ms) & 15.6 & 14.9 & -4.7\% \\
Layers 4 GiB, batch 2048 (ms) & 351.6 & 339.0 & -3.6\% \\
\bottomrule
\end{tabular}

\caption{Working sets that fit in VRAM: managed memory performs like device memory.}
\label{tab:overhead}
\end{table}

Table~\ref{tab:overhead} compares unmodified PyTorch with the same program under lowvram when
everything fits. GEMM throughput and sweep bandwidth agree within 1\%; the layer stack differs by
a few percent in both directions, within the spread of the individual repetitions (the ranges
overlap; a laptop GPU's clocks vary with temperature). After the first touch, managed pages are
resident and accessed at full device bandwidth. lowvram can therefore be left on: it costs nothing
until it is needed.

\subsection{Cyclic access beyond VRAM}

\begin{figure}[h]
\centering
\fig{0.72\linewidth}{sweep}
\caption{Effective bandwidth of a working set touched once per pass in a fixed order.}
\label{fig:sweep}
\end{figure}

Figure~\ref{fig:sweep} shows the cost of oversubscription in its purest form. While the working
set fits, both configurations run at about \num{sweepFitGBps}\,GB/s; unmodified PyTorch fails
beyond \num{usableVram}\,GiB, and lowvram drops to \num{sweepOversubGBps}\,GB/s at every larger
size. Cyclic access over a set larger than memory is the worst case for least-recently-used
eviction: the page needed next is always the one evicted longest ago, so every access misses and
each pass migrates the whole set over PCIe. The effective rate sits slightly above the pinned-copy
bandwidth because migration runs in both directions at once (pages in, evicted pages out) and the
metric counts bytes read plus written. The cliff is the whole story of demand paging: no
gradual degradation, but full speed or PCIe speed, depending on whether the set that is
\emph{reused} fits.

\subsection{Layer stack: paging versus streaming, and what hints recover}

\begin{figure}[h]
\centering
\fig{\linewidth}{layers}
\caption{Seconds per forward pass of the layer stack. Unmodified PyTorch fails beyond the free
VRAM (shaded).}
\label{fig:layers}
\end{figure}

\begin{table}[h]
\centering\small
\begin{tabular}{lrrrrrr}
\toprule
 & \multicolumn{3}{c}{batch 16} & \multicolumn{3}{c}{batch 2048} \\
\cmidrule(lr){2-4}\cmidrule(lr){5-7}
Mode & 6\,GiB & 8\,GiB & 12\,GiB & 6\,GiB & 8\,GiB & 12\,GiB \\
\midrule
Unmodified & OOM & OOM & OOM & OOM & OOM & OOM \\
lowvram & 1.83 & 2.49 & 3.76 & 2.20 & 3.01 & 4.72 \\
lowvram + ReadMostly & 1.26 & 1.61 & 2.70 & 1.63 & 2.32 & 3.30 \\
lowvram + prefetch & 2.98 & 3.75 & 5.92 & 3.34 & 4.50 & 6.64 \\
lowvram + remote map & 0.33 & 0.51 & 1.12 & 4.57 & 9.12 & 16.98 \\
manual streaming & 1.06 & 1.46 & 2.25 & 1.57 & 2.08 & 3.10 \\
manual streaming, overlapped & 1.05 & 1.38 & 2.24 & 1.05 & 1.41 & 2.07 \\
CPU (FP32) & -- & 0.96 & -- & -- & 27.85 & -- \\
\bottomrule
\end{tabular}

\caption{Layer stack beyond VRAM, seconds per forward pass (median).}
\label{tab:layers}
\end{table}

The layer stack reads every weight once per pass, in order: the same cyclic pattern, now with
computation in between (Figure~\ref{fig:layers}, Table~\ref{tab:layers}; numbers below are for
8\,GiB of weights).

\paragraph{Paging versus streaming.} lowvram alone takes \num{layersLoLv}\,s per pass at batch 16
and \num{layersHiLv}\,s at batch 2048, against \num{layersLoOff} and \num{layersHiOff}\,s for
manual streaming and \num{layersLoOvl} and \num{layersHiOvl}\,s for overlapped streaming. Demand
paging is thus 1.4--2.1$\times$ slower than a hand-written scheme. Three properties explain the
gap: transfers start only after a fault, in small granules, instead of as large DMA copies; they
do not overlap computation, because the kernel that needs the data is the one waiting for it; and
eviction copies every page back to host memory, although the weights never change, so PCIe
carries each byte twice.

\paragraph{What one hint recovers.} The hint modes test how much of that gap a program could
close by telling the driver one thing it knows. \texttt{READ\_MOSTLY} on the weights removes the
write-back: evicting a read-duplicated page drops it. That alone brings paging to \num{layersLoRm}
and \num{layersHiRm}\,s, about 10\% slower than manual streaming at batch 16. Prefetching the next two
layers makes things \emph{worse} (\num{layersLoPf} and \num{layersHiPf}\,s): when memory is full,
every prefetch forces an eviction, and the prefetched pages compete with the ones in use.
Mapping the layers beyond a 3.5\,GiB budget for direct access from host memory, so that they never
migrate, is the fastest configuration of all for the memory-bound case (\num{layersLoRem}\,s,
2.7$\times$ faster than overlapped streaming): the first 3.5\,GiB stay resident, and the rest are
read once over PCIe, at the link's speed, without copies or faults. The same mapping is the
slowest for the compute-bound case (\num{layersHiRem}\,s), because a large matrix product reads
each weight many times, and every one of those reads crosses PCIe.

\paragraph{CPU crossover.} At batch 16 the CPU computes the 8\,GiB stack in \num{layersLoCpu}\,s,
faster than lowvram, manual streaming and overlapped streaming, although it reads FP32 weights,
twice the bytes: it reads host memory directly at tens of GB/s, while every GPU mode is limited by
the 6\,GB/s link. At batch 2048 the CPU needs \num{layersHiCpu}\,s, and every GPU mode wins by an
order of magnitude.

The lesson for lowvram is that the mechanism is sound but uninformed. Its default---migrate on
fault, evict by recency, write back---is the right one for data that is reused and modified, and
the wrong one for read-only data streamed once per pass. A hint that only the program can give
recovers most of the difference.

\subsection{A real model: YuE2-3B}

\begin{figure}[h]
\centering
\fig{\linewidth}{yue2}
\caption{YuE2-3B on a 6\,GiB GPU. Left: token generation, which uses half of the weights.
Right: a flow-matching step, which uses all of them.}
\label{fig:yue2}
\end{figure}

\begin{table}[h]
\centering\small
\begin{tabular}{lrrrr}
\toprule
 & \multicolumn{2}{c}{Token generation (AR)} & \multicolumn{2}{c}{Flow-matching step (NAR)} \\
\cmidrule(lr){2-3}\cmidrule(lr){4-5}
Mode & prompt tok/s & decode tok/s & s/step & GPU weights \\
\midrule
Unmodified & OOM & OOM & OOM & all \\
Accelerate offload & 607 & 1.9 & 1.32 & part \\
CPU (BF16) & 46 & 6.8 & 52.18 & none \\
lowvram & 611 & 43.1 & 1.90 & all (paged) \\
AR half only (oracle) & 1694 & 45.3 & -- & AR half \\
\bottomrule
\end{tabular}

\caption{YuE2-3B, median over \num{yueRepeats} runs (CPU: one run).}
\label{tab:yue2}
\end{table}

Unmodified PyTorch cannot load YuE2 on this GPU. Under lowvram the same code generates
\num{yueLvDecode} tokens per second (range over runs \num{yueLvDecodeRange}): \num{yueLvOfOracle}\%
of the \num{yueOracleDecode} tokens per second of the oracle, which runs the AR half from ordinary
device memory after we deleted the NAR half by hand. Accelerate, the standard way to run an
oversized Hugging Face model, reaches \num{yueAccDecode} tokens per second---lowvram is
\num{yueLvOverAcc}$\times$ faster---and the CPU \num{yueCpuDecode}.

The difference is placement granularity. Accelerate places whole decoder layers, the unit the
model declares as unsplittable, and each layer on the CPU side carries its NAR attention and MLP
along with the AR ones; for every token it copies all offloaded layers to the GPU, including the
half that token generation never touches. Demand paging needs no declaration: the NAR pages are
never accessed during generation, so they are evicted once and stay in host memory, and the
3.4\,GiB the AR path reads every token stay resident. Figure~\ref{fig:warmup} shows the transient:
the first tokens fault the AR weights in, then the per-token latency settles within a few percent of
the oracle's. Prompt processing, which runs once, pays for that first touch: \num{yueLvPrefill}
tokens per second against \num{yueOraclePrefill} for the oracle and \num{yueAccPrefill} for
Accelerate.

The flow-matching step is the opposite case. It runs both halves over a 1,538-token sequence, so
all 6.8\,GiB are touched every step, in layer order: the cyclic pattern of
Figure~\ref{fig:sweep}. lowvram needs \num{yueLvNar}\,s per step against \num{yueAccNar}\,s for
Accelerate, which keeps 4\,GiB resident and streams the rest in planned copies; paging is
\num{yueNarSlowdown}$\times$ slower. The CPU needs \num{yueCpuNar}\,s per step.

The same program therefore shows both regimes. When the data a phase reuses fits, transparent
paging finds it at page granularity, without being told, and runs close to native speed; when it
does not, a planned scheme with knowledge of the model wins.

\subsection{A dense model: Qwen3-4B}

\begin{figure}[h]
\centering
\fig{\linewidth}{llm}
\caption{Qwen3-4B (\num{llmWeightsGiB}\,GiB of BF16 weights) on the same GPU.}
\label{fig:llm}
\end{figure}

YuE2's split hot set favors demand paging. A dense language model is the general case and the
unfavorable one: every generated token reads every weight, so the reused set is the whole model.
Qwen3-4B holds \num{llmWeightsGiB}\,GiB of BF16 weights (Figure~\ref{fig:llm},
Table~\ref{tab:llm}). Unmodified PyTorch cannot load it. Under lowvram, token generation runs at
\num{llmLvDecode} tokens per second: each token migrates the model through the 6\,GB/s link, in
both directions, as in the sweep. Accelerate reaches \num{llmAccDecode} tokens per second by keeping
4\,GiB resident and copying only the rest, once, in planned transfers; the CPU reaches
\num{llmCpuDecode}, the fastest of the three, because it reads host memory directly instead of
through the link. Prompt processing reuses each weight across the prompt's 1,024 tokens, so
transfers matter less: lowvram reaches \num{llmLvPrefill} tokens per second, against
\num{llmAccPrefill} for Accelerate and \num{llmCpuPrefill} for the CPU. Here lowvram is the
slowest GPU option but still an order of magnitude faster than the CPU.

\begin{table}[h]
\centering\small
\IfFileExists{generated/table_llm.tex}{\begin{tabular}{lrr}
\toprule
Mode & prompt tok/s & decode tok/s \\
\midrule
Unmodified & OOM & OOM \\
Accelerate offload & 837 & 1.17 \\
CPU (BF16) & 25 & 3.63 \\
lowvram & 499 & 0.58 \\
\bottomrule
\end{tabular}
}{\na}
\caption{Qwen3-4B, median over \num{yueRepeats} runs (CPU: one run).}
\label{tab:llm}
\end{table}

For dense models larger than VRAM, demand paging is the wrong tool: a planned split is
faster in both phases, and for token generation a scheme that computes the remainder on the
CPU---as hybrid engines such as llama.cpp do---is faster still. lowvram turns the failure into a
run, but when a framework offers offloading, it should be preferred for such models.

\begin{figure}[h]
\centering
\fig{0.72\linewidth}{yue2_warmup}
\caption{Per-token latency of YuE2 token generation under lowvram and the oracle.}
\label{fig:warmup}
\end{figure}

% ---------------------------------------------------------------------------
\section{Where lowvram is valuable}
\label{sec:where}

The measurements reduce to one question: \emph{does the data that a phase reuses fit in
VRAM?} The total size of the program does not matter; the reused set does.

\begin{table}[h]
\centering\small
\begin{tabular}{p{4.6cm}p{2.6cm}p{6.6cm}}
\toprule
Situation & lowvram & Evidence and alternative \\
\midrule
Program fits & free & Table~\ref{tab:overhead}; can be left on \\
Total exceeds VRAM, reused set fits (multi-path models, pipelines of several models, rarely used
modules) & \textbf{best option} & YuE2 generation: \num{yueLvOfOracle}\% of native,
\num{yueLvOverAcc}$\times$ Accelerate; no alternative finds the hot set without being told \\
Transient peak (one large batch, a decode step at high resolution) & \textbf{best option} & a
crash becomes a slow step; offloading the whole run to avoid one peak costs every step \\
Reused set exceeds VRAM, compute-bound & works, 1.4--2$\times$ slower than planned streaming &
YuE2 flow step, layer stack at batch 2048; use framework offloading if available \\
Reused set exceeds VRAM, memory-bound (dense LLM decoding) & works, but slowest & Qwen3-4B,
layer stack at batch 16; CPU or a hybrid engine (llama.cpp) is faster \\
Multi-process GPU sharing, \texttt{cudaMallocAsync}, \texttt{cuMemCreate} & not covered & allocations
bypass the hook \\
\bottomrule
\end{tabular}
\caption{When transparent oversubscription is the right tool.}
\label{tab:where}
\end{table}

In practice the comparison is often not with an optimized offload engine but with a program that
cannot run at all, or with porting it to a framework that can offload. lowvram makes the program
run first; Table~\ref{tab:where} tells whether it is then worth the effort to do better.

\paragraph{What could close the gap.} The hint experiments point at a middle ground between a
transparent shim and a framework. A program---or a framework integration---that marks its weights
\texttt{READ\_MOSTLY} recovers most of the distance to hand-written streaming, and one that knows
which weights are read once per pass can map them for direct access and beat streaming in the
memory-bound regime. Neither hint is safe to apply blindly to every allocation: read-mostly pages
are expensive to write, and direct access is slow for data that is reused. Applying them from
allocation-time heuristics (size, first-write behavior) is future work.

% ---------------------------------------------------------------------------
\section{Limitations and threats to validity}

One machine, one GPU generation and one PCIe configuration (x8); a desktop GPU with an x16 link
has twice the transfer bandwidth and would shift every oversubscribed number. The GPU also drove
the display. YuE2 inputs are synthetic; the sampling loop is greedy; we do not measure audio
quality, which lowvram cannot change. Reported free memory is inflated by a constant, which can
lead a program to exhaust host RAM. The comparison with Accelerate uses one weight budget.
\todo{CPU baseline in BF16 on a CPU without native BF16; llama.cpp hybrid not measured.}

% ---------------------------------------------------------------------------
\section{Related work}

\paragraph{Unified Memory.} Unified Memory oversubscription and its performance have been studied
extensively: Landaverde et al.\ \cite{landaverde2014} measured early Unified Memory performance;
Zheng et al.\ \cite{zheng2016} examined paged GPU memory and the role of prefetching; Ganguly et
al.\ \cite{ganguly2019} the interplay of prefetching and eviction under oversubscription; DeepUM
\cite{deepum} adds correlation prefetching for DNN training. TensorFlow exposed oversubscription
through a memory fraction above one \cite{tf-unified}, and the paged optimizers of
bitsandbytes \cite{qlora} place optimizer state in managed memory. lowvram applies the mechanism to
unmodified programs, below any framework.

\paragraph{API interposition.} Intercepting the CUDA API is how GPUs are remoted and shared:
rCUDA \cite{rcuda} and GVirtuS \cite{gvirtus} forward calls to remote GPUs, and Kubernetes
sharing layers such as HAMi \cite{hami} interpose on allocation to enforce quotas. lowvram uses the
same technique for a different end, and documents the lookup path a modern CUDA runtime requires.

\paragraph{Tensor swapping and offloading.} vDNN \cite{vdnn} offloads activations during training;
SwapAdvisor \cite{swapadvisor} and Capuchin \cite{capuchin} plan swapping and recomputation;
ZeRO-Offload and ZeRO-Infinity \cite{zero-offload,zero-infinity} offload optimizer state and
parameters to CPU and NVMe; FlexGen \cite{flexgen} schedules LLM inference across GPU, CPU and
disk; Accelerate \cite{accelerate} places modules across devices. These systems use knowledge of
the computation to outperform demand paging, but require running within their framework.

\paragraph{Hybrid CPU--GPU inference.} llama.cpp \cite{llamacpp} computes the layers that do not
fit on the CPU; PowerInfer \cite{powerinfer} splits hot and cold neurons between GPU and CPU;
KTransformers \cite{ktransformers} and expert offloading \cite{eliseev2023} do so for
mixture-of-experts models. For memory-bound decoding they beat any scheme that moves weights over
PCIe, at the cost of model-specific engines.

% ---------------------------------------------------------------------------
\section{Conclusion}

A small interposition library---three replaced driver functions plus the lookups that reach
them---runs unmodified CUDA programs beyond device memory. It costs nothing when a program fits and
turns out-of-memory failures into runs when it does not. Its speed beyond VRAM is decided by one
property of the workload, whether the data a phase reuses fits: when it does, demand paging finds
it at page granularity without being told, and a 7.3\,GB music model runs on a 6\,GiB laptop GPU
at \num{yueLvOfOracle}\% of native speed, \num{yueLvOverAcc}$\times$ faster than standard
offloading; when it does not, paging is bound by PCIe and a planned or CPU-based scheme is faster.
A single driver hint from the program recovers most of that difference, which marks the next step:
moving from a transparent shim towards an informed one, without giving up the property that makes
it useful---no change to the program.

\bibliographystyle{plain}
\bibliography{references}
\end{document}
